\section{Difficulties and Solutions}
\label{sec:difficulties}

\subsection{Data Leakage in a Time Series}
\label{sec:leakage}

A random row split mixes later hours into the training set. Neighbours of a
test hour then appear on both sides of the split, and lagged copies of the
target leak the answer. The resulting MAE looks strong and does not survive
the next week of operations.

The project uses a blocked temporal split. \(D_{\mathrm{train}}\),
\(T_{\mathrm{valid}}\) and \(T_{\mathrm{test}}\) are contiguous, ordered and
disjoint (Section~\ref{sec:models}). Rolling means and rolling standard
deviations are computed on the demand series after a one-hour shift, so hour
\(t\) is absent from the features that predict it. Validation labels are used
to choose tree size, learning rate and the ridge penalty. They are not used
as training targets, and the test block is not used to pick a model.

\subsection{Non-linear Weather Response}
\label{sec:nonlinear}

A straight line in temperature understates the load on hot afternoons, when
cooling equipment switches on and demand bends upward. Humidity and radiation
move with temperature, so a single coefficient is a poor summary of the
weather effect.

Two complementary fixes are in the feature module and in the model module.
The feature module adds cooling degree, squared temperature and the
temperature-by-hour product, which give the linear model a limited ability to
bend. The tree models (random forest, XGBoost, LightGBM) partition
temperature and hour into regions, so the same hot hour can receive a
different leaf value from a mild hour. The scatter plot in
Figure~\ref{fig:weather} is the diagnostic that this structure is worth
including. The gain shares in Figure~\ref{fig:importance} are the check that
the fitted model actually uses it.

\subsection{Error at the Peak, Hidden by the Mean}
\label{sec:peak-error}

MAE averages every hour. A quiet night with a small error can offset an
evening miss that crosses the capacity threshold. An operator who sees only
MAE will not know whether the peak hour and the peak level were forecast in
time.

The evaluation module therefore scores peaks separately
(Section~\ref{sec:peak}). Peak magnitude error measures the level miss.
Peak timing error measures whether the forecast peak falls on the right hour.
Breach anticipation rate measures whether hours above \(\tau\) were flagged
before they occurred. The dashboard repeats the same threshold on the plotted
forecast, so the numeric BAR and the on-screen alert refer to one rule.
